% =====================================================================
%  4.5  Atividade 3.2.2.5 -- Esgotamento descentralizado na área do IDR
% =====================================================================
\subsection{Atividade 3.2.2.5 -- Instalar sistemas descentralizados de esgotamento sanitário na área de atuação do IDR-Paraná}
\label{subsec:esgotamento-idr}

\subsubsection{Objetivo e fundamentação}

Instalar 2.300 sistemas descentralizados de esgotamento sanitário em
propriedades rurais situadas nas áreas prioritárias do PSH-PR de atuação do
IDR-Paraná, preferencialmente em mananciais de abastecimento público
(\autoref{subsec:criterios-idr}). O processo terá início após a elaboração dos
Planos Integrados de Propriedade (PIP) pelo IDR-Paraná \cite{mop2026}.

\begin{fichaatividade}{Ficha-síntese da Atividade 3.2.2.5}{qua:ficha-3225}
  \campo{Código}{3.2.2.5}
  \campo{Tipo de atividade}{Serviço e obra}
  \campo{Forma de execução}{Contratação pública (licitação)}
  \campo{Fonte de recursos}{BIRD}
  \campo{Responsável}{IAT / IDR-Paraná}
  \campo{Dependências institucionais}{--}
  \campo{Prazo estimado}{Ano 2 ao ano 5}
  \campo{Produto esperado}{2.300 sistemas implantados}
  \campo{Unidade de medida}{Sistema implantado}
  \campo{Evidência de comprovação}{Relatório técnico aprovado e comprovante de pagamento}
  \campo{Periodicidade}{Por lote}
  \campo{Validação}{IAT / IDR-Paraná}
\end{fichaatividade}

\subsubsection{Procedimentos}

Os procedimentos são os mesmos da \atividade{3.2.2.4}
(\autoref{subsubsec:proc-esgotamento} e \autoref{fig:fluxograma-esgotamento}),
da contratação à finalização. As diferenças entre as duas atividades estão
no \autoref{qua:diferencas-esgotamento}.

\begin{quadro}[htbp]
\caption{Diferenças entre as Atividades 3.2.2.4 e 3.2.2.5}
\label{qua:diferencas-esgotamento}
\small
\renewcommand{\arraystretch}{1.25}
\begin{tabularx}{\textwidth}{|P{3.4cm}|L|L|}
\hline
\cab{Aspecto} & \cab{3.2.2.4 -- Comunidades rurais} & \cab{3.2.2.5 -- Área do IDR-Paraná} \\ \hline
Origem da demanda & Diagnóstico do trabalho técnico-social (\atividade{3.2.2.2}) &
  Planos Integrados de Propriedade (PIP) elaborados pelo IDR-Paraná \\ \hline
Área de atuação & Comunidades atendidas pelo abastecimento de água (\atividade{3.2.2.3}) &
  Mananciais prioritários do PSH-PR na área de atuação do IDR-Paraná \\ \hline
Critérios de seleção & \autoref{subsec:criterios-comunidades} & \autoref{subsec:criterios-idr} \\ \hline
Projeto & Projeto individual por propriedade & Projeto sanitário doméstico da propriedade \\ \hline
Meta & 4.067 sistemas (70\% dos domicílios com abastecimento) & 2.300 sistemas \\ \hline
Dependências & Prefeituras & -- \\ \hline
Validação & IAT & IAT / IDR-Paraná \\ \hline
\end{tabularx}
\fonte{Adaptado de \citeonline{mop2026}.}
\end{quadro}

\pendencia{No MOP, o texto das \atividade{3.2.2.4} e \atividade{3.2.2.5} é
praticamente idêntico. Nesta estrutura, o procedimento é descrito uma vez e
as diferenças vão para o \autoref{qua:diferencas-esgotamento}. Confirmar se o
fluxo de validação (auditoria, verificação \textit{in loco}) é o mesmo nas
duas.}

\subsubsection{Interfaces}

\begin{itemize}
  \item Subcomponente 2.3 do PSH-PR: elaboração dos PIP pelo IDR-Paraná.
  \item \atividade{3.2.2.4}: procedimento comum de contratação, instalação,
        verificação e pagamento.
\end{itemize}
